\renewcommand{\Anno}{1996}

\renewcommand{\NumeroProblema}{1}

\renewcommand{\Testo}{
Siano $x$ e $y$ numeri reali positivi tali che $x+y<\pi/2$. Dimostrare che
$$\sin (x+y) < \sin x + \sin y < \tan (x+y).$$}

\renewcommand{\Risposta}{
% SSSUP-BOZZA-AI: Risposta; autore=ChatGPT; data=2026-11-08; stato=DA-VERIFICARE
% SSSUP-PROVENIENZA: AI: elaborazione nuova dal testo
Le due disuguaglianze sono strette. Si possono ottenere dalla catena
\[
\sin(x+y)<\sin x+\sin y<\tan x+\tan y<\tan(x+y).
\]
}

\renewcommand{\Aiutino}{Vi sono almeno due approcci possibili:
\begin{itemize}
	\item  uno algebrico, basato sulle formule di addizione del seno e 
	della tangente, e sulle ovvie disuguaglianze $\cos x<1$ e $\sin x<\tan 
	x$ (e analoghe per $y$, ovviamente nella zona in questione);

	\item  uno geometrico, basato sull'interpretazione del seno e della tangente sulla circonferenza trigonometrica.
\end{itemize}
}

\renewcommand{\SoluzioneUfficiale}{
Dalla formula di addizione del seno abbiamo che
$$\sin(x+y)=\sin x\cos y+\sin y\cos x<\sin x+\sin y,$$
dove l'ultima maggiorazione segue dal fatto che i due seni sono positivi ed i due coseni sono minori di 1 (per le ipotesi fatte $x$ e $y$ appartengono all'intervallo $(0,\pi/2)$).  Questo dimostra la prima disuguaglianza richiesta.

Dalla formula di addizione della tangente abbiamo invece che
$$\tan(x+y)=\frac{\tan x+\tan y}{1-\tan x\tan y}>
\frac{\tan x+\tan y}{1}>\sin x+\sin y,$$
dove la prima minorazione segue dal fatto che $0<1-\tan x\tan y<1$ (questo perchè nelle ipotesi date $\tan x$, $\tan y$ e $\tan(x+y)$ sono quantità positive), e la seconda segue dal fatto che le tangenti sono maggiori dei rispettivi seni (e qui si usa ancora una volta che $x$ e $y$ appartengono all'intervallo $(0,\pi/2)$).  Questo
dimostra la seconda disuguaglianza richiesta.  }

\renewcommand{\SoluzioniAlternative}{\TitoloSoluzioneAlternativaA 
Con riferimento alla Figura 1, sia $OA=1$, sia $x$ la misura dell'arco $AP$ e $y$ la misura dell'arco $PQ$: allora
$$QL=\sin(x+y),
\hspace{3em}
QK=\sin y,
\hspace{3em}
NL=PH=\sin x.$$

D'altra parte $QM<QK$ (in quanto, rispettivamente, cateto ed ipotenusa nel triangolo rettangolo $QMK$) e quindi
$$\sin(x+y)=QL<QM+NL<QK+PH=\sin y+\sin x.$$

Questo dimostra la prima disuguaglianza.

\hfill
\psset{unit=10mm}
\pspicture(-1,-2)(4.5,7)
\rput{0}(1.75,-1.5){Figura 1}
\pstGeonode[PosAngle={315,315,25,55}](0,0){O}(4;0){A}(4;25){P}(4;55){Q}
\pstProjection{O}{A}{P}[H]
\pstProjection{O}{A}{Q}[L]
\pstProjection{O}{P}{Q}[K]
\pstProjection[PosAngle=160]{Q}{L}{P}[N]
\pstProjection[PosAngle=200]{Q}{L}{K}[M]
\rput{0}(4.4;13){$x$}
\rput{0}(4.4;40){$y$}

\psline{->}(-1,0)(4.5,0)
\psline{->}(0,-1)(0,4.5)
\psarc{-}{4}{-10}{100}

\pstLineAB{O}{P}
\pstLineAB{Q}{O}

\psset{linecolor=blue,linewidth=.7\pslinewidth}
\pstLineAB{Q}{L}
\pstLineAB{P}{H}
\pstLineAB{Q}{K}
\pstRightAngle[RightAngleSize=0.25]{Q}{K}{O}

\pstLineAB[linestyle=dashed]{N}{P}
\pstLineAB[linestyle=dashed]{M}{K}

\endpspicture
\hspace{4em}
\psset{unit=10mm}
\pspicture(-1,-2)(5.5,7)
\rput{0}(2.25,-1.5){Figura 2}
\pstGeonode[PosAngle={315,315,80,90}](0,0){O}(4;0){A}(4;25){P}(4;55){Q}
\pstProjection{O}{A}{P}[H]
\pstProjection{O}{P}{Q}[K]
\pstRotation[RotAngle=-90,PointSymbol=none]{A}{O}[A']
\pstInterLL[PosAngle=-20]{A}{A'}{O}{P}{B}
\pstInterLL{A}{A'}{O}{Q}{C}
\pstProjection{O}{B}{C}[D]

\psline{->}(-1,0)(4.5,0)
\psline{->}(0,-1)(0,4.5)
\psarc{-}{4}{-10}{100}

\pstLineAB{O}{Q}
\pstLineAB{O}{P}

\psset{linecolor=blue,linewidth=.7\pslinewidth}
\pstLineAB{P}{H}
\pstLineAB{Q}{K}
\pstLineAB{Q}{C}
\pstLineAB{D}{C}
\pstLineAB{D}{P}
\pstLineAB[nodesepB=-0.5]{A}{C}

\pstRightAngle[RightAngleSize=0.25]{Q}{K}{O}
\pstRightAngle[RightAngleSize=0.25]{C}{D}{O}

\endpspicture
\hfill\mbox{}

Nella Figura 2, in cui gli archi $AP$ e $PQ$ misurano ancora, rispettivamente, $x$ e $y$, abbiamo invece che
\begin{itemize}
	\item  $PH<BA$ (per similitudine tra i triangoli $OPH$ e $OBA$);

	\item  $QK<CD$ (per similitudine tra i triangoli $OQK$ e $OCD$);

	\item $CD<CB$ (in quanto, rispettivamente, cateto ed ipotenusa nel triangolo rettangolo $ODC$).
\end{itemize}
Ne segue che
$$\sin x+\sin y=QK+PH<CD+BA<CB+BA=CA=\tan(x+y),$$
il che dimostra la seconda disuguaglianza.
}

\renewcommand{\Osservazioni}{}

\renewcommand{\Commenti}{}

